\sgaStableTarget{sga6.e13.prop.4-5}\textbf{Proposition 4.5. (Functorial behavior).} \emph{Let @@M0@@ be a morphism between two proper @@M1@@-schemes, and let @@M2@@ be an invertible module on @@M3@@. Then}~:
\begin{enumerate}[label=(\roman*)]
\item \emph{If @@M4@@ is @@M5@@-equivalent to @@M6@@ (resp. numerically equivalent to @@M7@@), then so is @@M8@@.}
\item \emph{Conversely, when @@M9@@ is surjective, if @@M10@@ is @@M11@@-equivalent to @@M12@@ (resp. numerically equivalent to @@M13@@), then so is @@M14@@.}
\end{enumerate}

Indeed, assertion (i) (resp. (ii)) for @@M15@@-equivalence follows at once from 4.1 (resp. 4.3). On the other hand, let @@M16@@ be an integral closed curve in @@M17@@, and @@M18@@ its image, endowed with the induced reduced structure. When @@M19@@, we have the projection formula (X 4.4.4)
\sgaStableTarget{sga6.e13.eq.4-5-1}@@M20@@
where @@M21@@ (resp. @@M22@@) is the function field of @@M23@@ (resp. @@M24@@), and when @@M25@@, we have @@M26@@. Consequently, if @@M27@@ is numerically equivalent to @@M28@@, then @@M29@@\,; hence (i).

Conversely, when @@M30@@ is surjective, given an integral closed curve @@M31@@ in @@M32@@, one can find an integral closed curve @@M33@@ in @@M34@@ such that @@M35@@: let @@M36@@ be the generic point of @@M37@@, @@M38@@ a point of the fiber @@M39@@ algebraic over @@M40@@, and @@M41@@ the closure of @@M42@@. Then 4.5.1 shows that @@M43@@ implies @@M44@@\,; hence (ii).

\sgaStableTarget{sga6.e13.thm.4-6}\textbf{Theorem 4.6. (The numerical characterization of @@M45@@).} \emph{Let @@M46@@ be an algebraically closed field, @@M47@@ a proper @@M48@@-scheme, and @@M49@@ an invertible sheaf on @@M50@@. Then the following conditions are equivalent}~:
\begin{enumerate}[label=(\roman*)]
\item \emph{@@M51@@ is @@M52@@-equivalent to @@M53@@.}
\item \emph{For every coherent sheaf @@M54@@ on @@M55@@, we have}
